\section{HD\,189733\,b}
\label{sec:hd189733b}

HD\,189733\,b is the brightest transiting hot Jupiter and has become the
archetypal target for upper-atmosphere and escape diagnostics. Its extended
thermosphere has been probed in H\,\textsc{i} Ly$\alpha$, the Balmer
H$\alpha$ line, the metastable He\,\textsc{i} 10830\,\AA{} triplet, and a set
of metal resonance lines (Na\,\textsc{i} D, O\,\textsc{i}, C\,\textsc{ii},
Si\,\textsc{iii}, N\,\textsc{v}, and the near-ultraviolet Mg/Fe region).
A recurring theme is the strong stellar activity of the K-dwarf host, which
contaminates the chromospheric-line diagnostics (H$\alpha$, Ca\,\textsc{ii},
Na\,\textsc{i} D) and drives night-to-night and epoch-to-epoch variability in
the reported signals.

\subsection{Observations}

\begin{longtable}{@{}p{0.14\textwidth}p{0.20\textwidth}p{0.44\textwidth}p{0.16\textwidth}@{}}
\toprule
Line(s) & Instrument & Measured result & Reference \\
\midrule
\endfirsthead
\toprule
Line(s) & Instrument & Measured result & Reference \\
\midrule
\endhead
\midrule
\multicolumn{4}{r}{\emph{continued on next page}}\\
\endfoot
\bottomrule
\endlastfoot

Ly$\alpha$ (H\,\textsc{i}); O\,\textsc{i}, C\,\textsc{ii} limits &
HST/ACS &
Ly$\alpha$ transit depth $5.05\pm0.75\%$ (3.5$\sigma$ above disk); upper limits O\,\textsc{i} $\lambda1305$ $<11.1\%$, C\,\textsc{ii} $\lambda1335$ $<5.5\%$; first Ly$\alpha$ escape signature &
\texttt{2010A\&A...514A..72L} \\

Ly$\alpha$ (H\,\textsc{i}) &
HST/STIS &
Spectrally resolved; \emph{temporal variability}: no absorption at epoch 1, $14.4\pm3.6\%$ at epoch 2 over $-230$ to $-140$ km\,s$^{-1}$; high velocities require stellar-wind interaction &
\texttt{2012A\&A...543L...4L} \\

Ly$\alpha$; Si\,\textsc{iii} $\lambda1206.5$, N\,\textsc{v} $\lambda1240$ &
HST/STIS &
Detailed analysis of the Sept.\ 2011 Ly$\alpha$ transit (false-positive prob.\ $3.6\%$); variations in Si\,\textsc{iii}/N\,\textsc{v} ascribed to a bow shock or stellar variability &
\texttt{2013A\&A...551A..63B} \\

O\,\textsc{i} (FUV) &
HST/COS (archival) &
Neutral-oxygen in-transit absorption $\sim6.4\pm1.8\%$; requires super-solar abundance and/or super-thermal broadening; flux variability noted &
\texttt{2013A\&A...553A..52B} \\

C\,\textsc{ii} $\lambda1335.7$, H\,\textsc{i} Ly$\alpha$, O\,\textsc{i} &
HST/COS + STIS &
Tentative but repeatable C\,\textsc{ii} absorption $5.2\pm1.4\%$; H\,\textsc{i} consistent with prior Ly$\alpha$; only marginal O\,\textsc{i} &
\texttt{2023AJ....166...89D} \\

NUV continuum; Mg\,\textsc{ii}/Mg\,\textsc{i}, Fe\,\textsc{ii} region &
HST/STIS (NUV) &
Broad excess absorption at $\sim2350$ and $\sim2600$\,\AA{} (Fe\,\textsc{ii} bands); Mg\,\textsc{ii} core absorption ruled out at 2--4$\sigma$; no constraint on Mg\,\textsc{i} &
\texttt{2023A\&A...671A.170C} \\

H$\alpha$ &
HET/HRS &
First H$\alpha$ exosphere detection: $(-8.72\pm1.48)\times10^{-4}$ over a 16\,\AA{} band; implies excitation temperature $T_{\rm exc}\approx2.6\times10^{4}$\,K &
\texttt{2012ApJ...751...86J} \\

H$\alpha$, Ca\,\textsc{ii} H\&K, Na\,\textsc{i} D &
(HARPS-derived) &
Excess H$\alpha$ $0.84\%$/$1.21\%$ (5.2/9.9$\sigma$), Na\,\textsc{i} D $0.41\%$ (6.5$\sigma$), Ca\,\textsc{ii} up to $2.14\%$ (9.8$\sigma$); argues signal is largely chromospheric/stellar &
\texttt{2016MNRAS.462.1012B} \\

H$\alpha$, Ca\,\textsc{ii} H\&K (Balmer pre-transit) &
optical high-res spectra &
Pre-transit Balmer signal detected but variable in strength/timing; partly correlated with Ca\,\textsc{ii} H (stellar activity); bow-shock geometry not reproduced &
\texttt{2016AJ....152...20C} \\

H$\alpha$ (Balmer) &
HARPS + Keck/HIRES &
Seven H$\alpha$ transits over 2006--2015; strong epoch-to-epoch variation, no clean Ca\,\textsc{ii}--H$\alpha$ trend; signal likely from extended atmosphere &
\texttt{2017AJ....153..217C} \\

H$\alpha$ (out-of-transit) &
high-res optical &
Near-transit H$\alpha$ core changes are uncommon in out-of-transit monitoring, strengthening the case for abnormal near-transit variability &
\texttt{2017AJ....153..185C} \\

H$\alpha$; Ca\,\textsc{ii} IRT &
TIGRE &
\emph{Non-detection}: $<6.2$\,m\AA{} pre-transit and $<5.6$\,m\AA{} in-transit (99\% conf.) after activity correction; challenges bow-shock/absorption claims &
\texttt{2018A\&A...619A..96K} \\

He\,\textsc{i} 10830 &
CARMENES &
$0.88\pm0.04\%$ mid-transit (net blueshift $-3.5\pm0.4$ km\,s$^{-1}$); compact He atmosphere $\sim0.2\,R_{\rm P}$, column $\sim4\times10^{12}$ cm$^{-2}$; activity pseudo-signals up to 80\% &
\texttt{2018A\&A...620A..97S} \\

He\,\textsc{i} 10830 &
GIARPS (GIANO-B + HARPS-N) &
$0.75\pm0.03\%$ (25$\sigma$) in strongest triplet component over 5 transits; night-to-night variability from stellar activity &
\texttt{2020A\&A...639A..49G} \\

He\,\textsc{i} 10830 &
SPIRou &
Detection confirmed in a homogeneous 11-planet survey; He presence correlates with stellar mass and EUV flux; p-winds interpretation &
\texttt{2023A\&A...677A.164A} \\

He\,\textsc{i} 10830 &
SPIRou &
Detection confirmed among 15 planets; new $\dot{M}$/temperature constraints via 1D Parker-wind + NLTE models &
\texttt{2024A\&A...688A.179M} \\

Na\,\textsc{i} D &
HET/HRS &
First ground-based Na detection: $(-67.2\pm20.7)\times10^{-5}$ in a narrow band ($>3\sigma$); $\sim3\times$ deeper than HD\,209458\,b &
\texttt{2008ApJ...673L..87R} \\

Na\,\textsc{i} D &
HST/STIS &
Na core absorption $0.09\pm0.01\%$ (9$\sigma$) in 5\,\AA; T--P profile shows upper-atmospheric heating (1280\,K wings, 2800\,K cores) &
\texttt{2012MNRAS.422.2477H} \\

Na\,\textsc{i} D &
HARPS &
Spectrally resolved doublet: contrasts $0.64\%$ (D2)/$0.40\%$ (D1), excess $0.32\pm0.03\%$ (10$\sigma$); $T\approx2600$--$3270$\,K at 9800--12700 km &
\texttt{2015A\&A...577A..62W} \\

Na\,\textsc{i} D &
HARPS &
Confirmed: D2 $0.64\pm0.07\%$, D1 $0.53\pm0.07\%$; effective photosphere $\sim1.13\,R_{\rm P}$, net blueshift; telluric-correction method comparison &
\texttt{2021MNRAS.502.4392L} \\

Na\,\textsc{i} D &
HARPS / HARPS-N (SLOPpy) &
Na doublet recovered as validation of the public SLOPpy pipeline; consistent with prior HD\,189733\,b results &
\texttt{2022A\&A...667A..19S} \\

Na\,\textsc{i}, Li\,\textsc{i} &
ESPRESSO &
Na doublet confirmed (with careful POLD/activity correction using ANTARESS + EvE); Na and Li signatures extracted &
\texttt{2025A\&A...700A..41M} \\

Na\,\textsc{i} D &
HARPS-N (GAPS) &
Significant Na absorption on most nights in a homogeneous 10-planet search; atmospheric height follows $T_{\rm eq}$--gravity trend; variability characterized &
\texttt{2025A\&A...693A.316S} \\

\end{longtable}

\noindent\emph{Note:} several optical/UV chromospheric-line results
(H$\alpha$, Ca\,\textsc{ii}, Na\,\textsc{i} D) carry large stellar-activity
systematics; the H$\alpha$ non-detection of \texttt{2018A\&A...619A..96K}
directly contradicts parts of the earlier H$\alpha$/bow-shock literature.

\subsection{Model studies}

Several observational papers above also carry forward models; those model
components are described here together with the dedicated theory papers.

\paragraph{Lecavelier des Etangs et al.\ 2010 --- \texttt{2010A\&A...514A..72L}.}
\emph{Code/method:} numerical simulation of escaping neutral hydrogen driven by
stellar Ly$\alpha$ radiation pressure (particle/dynamical model).
\emph{Initial conditions:} atomic-hydrogen exosphere launched from the planet;
free parameters are the escape rate and the ionizing EUV flux.
\emph{Key results:} the Ly$\alpha$ light curve is fit with
$\dot{M}\approx10^{9}$--$10^{11}$\,g\,s$^{-1}$ and $F_{\rm EUV}=2$--40 times
solar; best fit $\dot{M}=10^{10}$\,g\,s$^{-1}$, $F_{\rm EUV}=20\times$ solar.

\paragraph{Bourrier \& Lecavelier des Etangs 2013 --- \texttt{2013A\&A...557A.124B}.}
\emph{Code/method:} 3D particle Monte-Carlo model of the escaping-atom
dynamics, forward-modeling Ly$\alpha$ absorption line profiles for both
HD\,209458\,b and HD\,189733\,b.
\emph{Initial conditions:} neutral H atoms released from the upper atmosphere,
subject to stellar radiation pressure, photoionization, and (for
HD\,189733\,b) charge exchange/collisions with stellar-wind protons; stellar
wind parameterized by temperature, velocity, density.
\emph{Key results:} for HD\,189733\,b the observed $-230$ to $-140$ km\,s$^{-1}$
velocities require stellar-wind interaction (radiation pressure alone
insufficient); constrains stellar wind $T\sim3\times10^{4}$\,K,
$v=200\pm20$ km\,s$^{-1}$, $n\sim10^{3}$--$10^{7}$ cm$^{-3}$, escape rate
$4\times10^{8}$--$10^{11}$\,g\,s$^{-1}$, ionizing flux 6--23 times solar;
reveals an ``escape-limited'' behavior.

\paragraph{Salz et al.\ 2016 --- \texttt{2016A\&A...586A..75S}.}
\emph{Code/method:} 1D spherically symmetric hydrodynamic simulations with the
PLUTO--CLOUDY interface (TPCI), coupling detailed photoionization/plasma
physics with an MHD solver; a grid of 18 strongly irradiated gas planets
including HD\,189733\,b.
\emph{Initial conditions:} high-energy (XUV) irradiated thermosphere; solar
composition; radiative heating/cooling including Ly$\alpha$ and free--free
emission resolved self-consistently.
\emph{Key results:} reproduces the Na\,\textsc{i} D-derived $T$--$P$ profile of
HD\,189733\,b; radiative cooling (Ly$\alpha$) is important for this planet;
compact/massive thermospheres (e.g.\ HAT-P-2\,b) found hydrodynamically stable;
predicts planetary Ly$\alpha$ absorption/emission signals.

\paragraph{Lamp\'on et al.\ 2021a --- \texttt{2021A\&A...647A.129L}.}
\emph{Code/method:} 1D hydrodynamic, spherically symmetric outflow model
coupled to a non-LTE model for the He\,\textsc{i} triplet population; applied to
HD\,189733\,b and GJ\,3470\,b.
\emph{Initial conditions:} isothermal-type transonic wind constrained by the
measured He\,10830 absorption; H density from Ly$\alpha$ used to further
constrain temperature, mass-loss rate, and H/He ratio.
\emph{Key results:} HD\,189733\,b has a compact, hot thermosphere with peak
$T=12400^{+400}_{-300}$\,K, very low mean molecular mass
(H/He $=99.2/0.8\pm0.1$), and near-full ionization above $1.1\,R_{\rm P}$;
small radial outflow velocities (contrasted with the extended GJ\,3470\,b).

\paragraph{Lamp\'on et al.\ 2021b --- \texttt{2021A\&A...648L...7L}.}
\emph{Code/method:} same 1D hydrodynamic + He-triplet NLTE framework, used to
estimate hydrogen recombination and advection rates from combined He\,10830 and
Ly$\alpha$ absorption.
\emph{Initial conditions:} as in 2021a, for HD\,209458\,b, HD\,189733\,b, and
GJ\,3470\,b.
\emph{Key results:} identifies the three theoretical escape regimes
observationally --- HD\,209458\,b is energy-limited, \textbf{HD\,189733\,b is
recombination-limited}, and GJ\,3470\,b is photon-limited; proposes these as
benchmark cases.

\paragraph{Dos Santos et al.\ 2023 --- \texttt{2023AJ....166...89D}.}
\emph{Code/method:} isothermal Parker-wind forward model of the outflow, used to
interpret HST UV absorption of C, O, and H nuclei.
\emph{Initial conditions:} outflow temperature $\sim12100$\,K, solar C
abundance; escaping-nuclei altitudes matched to the observations.
\emph{Key results:} outflow is mostly neutral within $\sim2\,R_{\rm P}$ and
singly ionized beyond; C\,\textsc{ii} absorption consistent with
$\dot{M}\approx1.1\times10^{11}$\,g\,s$^{-1}$; predicted O\,\textsc{i} depth is
only comparable to the measurement uncertainty; Ly$\alpha$ comparison implies
interaction with the stellar environment.

\paragraph{He\,10830 survey models --- \texttt{2023A\&A...677A.164A} (Allart et al.\ 2023) and \texttt{2024A\&A...688A.179M} (Masson et al.\ 2024).}
\emph{Code/method:} the SPIRou He\,10830 detections are interpreted with 1D
isothermal Parker-wind + NLTE radiative-transfer models (the open-source
\texttt{p-winds} code in Allart et al.; two independent 1D Parker/NLTE models in
Masson et al.).
\emph{Initial conditions:} isothermal transonic wind parameterized by
thermosphere temperature and mass-loss rate, with line-of-sight velocity and
thermosphere altitude fit to the resolved triplet.
\emph{Key results:} confirm the HD\,189733\,b He detection and provide
$\dot{M}$/temperature constraints; both stress tension with some prior
literature values and the difficulty of tying He presence uniquely to
$\dot{M}$.

\paragraph{Cubillos et al.\ 2023 --- \texttt{2023A\&A...671A.170C}.}
\emph{Code/method (model comparison):} solar-abundance hydrodynamic
upper-atmosphere models compared against the observed NUV excess.
\emph{Key results:} such models \emph{cannot} reproduce the amplitude of the
$\sim2350$/$\sim2600$\,\AA{} features with Fe\,\textsc{ii} absorption at solar
abundance, leaving the NUV excess unexplained.
